% ============================================================================
% METHODS BRIEF: DRUG RECOMMENDATION SYSTEM
% Publication Format
% ============================================================================

\documentclass[11pt,a4paper]{article}
\usepackage[utf8]{inputenc}
\usepackage{amsmath,amssymb}
\usepackage{booktabs}
\usepackage{geometry}
\usepackage{graphicx}
\usepackage{hyperref}
\usepackage{float}
\geometry{margin=1in}

\begin{document}

\section*{Multi-Objective Drug Recommendation System Methods}

\subsection*{Overview}

The recommendation system employs a multi-objective optimization framework to identify therapeutic alternatives that balance clinical efficacy with safety. For each candidate alternative drug $d_{\text{alt}}$, we compute a recommendation score:

\begin{equation}
\text{RecScore}(d_{\text{alt}}) = 0.40 \cdot T(d_{\text{alt}}) + 0.35 \cdot SAF(d_{\text{alt}}) + 0.25 \cdot R(d_{\text{alt}})
\end{equation}

\subsection*{Score Components}

\subsubsection*{Therapeutic Similarity Score $T(d_{\text{alt}})$ (Weight: 0.40)}

Evaluates whether the alternative can serve the same clinical purpose as the original drug:

\begin{itemize}
    \item \textbf{ATC Code Matching}: Pharmacological class equivalence based on Anatomical Therapeutic Chemical classification
    \item \textbf{Shared Protein Targets}: Similar mechanisms of action through common molecular targets
    \item \textbf{Disease Indication Overlap}: Therapeutic applicability confirmation
    \item \textbf{Metabolic Pathway Similarity}: Comparable pharmacokinetic profiles
\end{itemize}

\subsubsection*{Safety Improvement Score $SAF(d_{\text{alt}})$ (Weight: 0.35)}

Quantifies the reduction in DDI risk when the alternative is used with other drugs in the current regimen:

\begin{itemize}
    \item Contraindicated interactions are penalized heavily (weight: 4.0)
    \item Major interactions are penalized moderately (weight: 3.0)
    \item Moderate interactions receive standard weight (weight: 2.0)
    \item Minor interactions receive minimal penalty (weight: 1.0)
\end{itemize}

\subsubsection*{Risk Reduction Score $R(d_{\text{alt}})$ (Weight: 0.25)}

Measures the net change in severe interactions and PRI score improvement:

\begin{equation}
R(d_{\text{alt}}) = \frac{\Delta_{\text{contraindicated}} + \Delta_{\text{major}} + \alpha \cdot \Delta_{\text{PRI}}}{\text{max\_possible\_reduction}}
\end{equation}

where $\alpha$ is a scaling factor for PRI improvement contribution.

\subsection*{Case Study: High-Risk Cardiovascular Polypharmacy}

We validated the methodology using a clinically realistic 5-drug cardiovascular regimen:

\begin{table}[H]
\centering
\caption{Original High-Risk Regimen}
\label{tab:original_regimen}
\begin{tabular}{ll}
\toprule
\textbf{Drug} & \textbf{Class} \\
\midrule
Warfarin & Anticoagulant \\
Amiodarone & Antiarrhythmic \\
Digoxin & Cardiac Glycoside \\
Quinidine & Antiarrhythmic \\
Propranolol & Beta-blocker \\
\bottomrule
\end{tabular}
\end{table}

Analysis revealed 10 severe drug-drug interactions: 2 contraindicated (Warfarin-Amiodarone causing increased bleeding risk, and Digoxin-Propranolol causing severe bradycardia) and 8 major interactions.

\subsection*{DOAC Alternatives Comparison}

\begin{table}[H]
\centering
\caption{Multi-Objective Recommendation Scores for DOAC Alternatives}
\label{tab:doac_alternatives}
\begin{tabular}{@{}lcccccc@{}}
\toprule
\textbf{Drug} & \textbf{RecScore} & \textbf{Therap.} & \textbf{Safety} & \textbf{Risk$\downarrow$} & \textbf{Contra} & \textbf{Major} \\
\midrule
\textbf{Dabigatran} & \textbf{0.801} & 0.850 & 0.782 & 0.750 & 0 & 1 \\
Fondaparinux & 0.524 & 0.850 & 0.347 & 0.250 & 0 & 3 \\
Apixaban & 0.385 & 0.850 & 0.129 & 0.000 & 0 & 4 \\
Rivaroxaban & 0.385 & 0.850 & 0.129 & 0.000 & 0 & 4 \\
Edoxaban & 0.385 & 0.850 & 0.129 & 0.000 & 0 & 4 \\
\bottomrule
\end{tabular}
\end{table}

\subsection*{Results Summary}

Substituting Warfarin with Dabigatran achieved:
\begin{itemize}
    \item Contraindicated interactions: 2 $\rightarrow$ 1 (50\% reduction)
    \item Major interactions: 8 $\rightarrow$ 6 (25\% reduction)
    \item Total severe interactions: 10 $\rightarrow$ 7 (30\% reduction)
    \item Average PRI: 0.569 $\rightarrow$ 0.504 (11.5\% improvement)
\end{itemize}

\section*{Figures}

\begin{figure}[H]
\centering
\includegraphics[width=0.9\textwidth]{plots_and_data/drug_alternatives_heatmap.png}
\caption{Drug alternatives heatmap showing recommendation scores for various therapeutic substitution options.}
\label{fig:alternatives_heatmap}
\end{figure}

\begin{figure}[H]
\centering
\includegraphics[width=0.9\textwidth]{plots_and_data/drug_substitution_network.png}
\caption{Drug substitution network visualization showing recommended alternatives and their therapeutic relationships.}
\label{fig:substitution_network}
\end{figure}

\begin{figure}[H]
\centering
\includegraphics[width=0.9\textwidth]{plots_and_data/network_safe_alternatives.png}
\caption{Safe alternatives network highlighting drugs with improved safety profiles compared to high-risk options.}
\label{fig:safe_alternatives}
\end{figure}

\begin{figure}[H]
\centering
\includegraphics[width=0.9\textwidth]{plots_and_data/class_alternatives_summary.png}
\caption{Class-level alternatives summary showing therapeutic substitution patterns across drug categories.}
\label{fig:class_alternatives}
\end{figure}

\section*{Conclusion}

The multi-objective recommendation framework provides a systematic approach for optimizing multi-drug regimens by balancing therapeutic similarity (40\%), safety improvement (35\%), and risk reduction (25\%). This evidence-based approach supports clinical decision-making in complex polypharmacy scenarios.

\bibliographystyle{plain}
\begin{thebibliography}{2}
\bibitem{wishart2018drugbank} Wishart DS, et al. DrugBank 5.0: a major update to the DrugBank database for 2018. \textit{Nucleic Acids Res}. 2018;46(D1):D1074-D1082.
\bibitem{xiong2022ddinter} Xiong G, et al. DDInter: an online drug-drug interaction database. \textit{Nucleic Acids Res}. 2022;50(D1):D1200-D1207.
\end{thebibliography}

\end{document}
